\documentclass[11pt,a4paper]{article}

\usepackage[utf8]{inputenc}
\usepackage[T1]{fontenc}
\usepackage{amsmath,amssymb,amsthm}
\usepackage[margin=1in]{geometry}
\usepackage{parskip}
\usepackage{hyperref}
\usepackage{booktabs}

\newtheorem{theorem}{Theorem}[section]
\newtheorem{lemma}[theorem]{Lemma}
\newtheorem{conjecture}[theorem]{Conjecture}
\theoremstyle{definition}
\newtheorem{definition}[theorem]{Definition}

\title{\bf Prime Density on Hilbert Planes Predicts Zeta Zero Ordering}
\author{Rick Wesson\\Support Intelligence, Inc.}
\date{August 2026}

\begin{document}

\maketitle

\begin{abstract}
We report a computational investigation into whether prime density on 3D
Hilbert curves encodes the zeta zero spectrum. When primes are mapped onto
a Hilbert curve and the cumulative excess density on each Z-plane is used
as the potential in a discrete Schr\"{o}dinger operator, the eigenvalues
correlate with the imaginary parts of the Riemann zeta zeros at
$|r| = 0.985$ for small systems ($N=64$ planes, $4{,}096$ primes).
However, the correlation \emph{decreases} as more primes are included,
falling to $|r| = 0.860$ with $16$ million primes at the same order.
This is the behavior of a small-$N$ artifact, not a convergent signal.
We conclude that prime density on Hilbert planes does \emph{not} provide
a first-principles construction of the Hilbert--P\'{o}lya operator.
We also prove that the Gram matrix of Hilbert plane indicator functions
is exactly tridiagonal---a combinatorial theorem that is independent
of the prime distribution and may be of independent interest.
\end{abstract}

\section{Introduction}

The Riemann Hypothesis (RH) states that all non-trivial zeros of
$\zeta(s)$ lie on the critical line $\Re(s) = 1/2$. The
Hilbert--P\'{o}lya conjecture proposes that the imaginary parts
$\gamma_k$ of these zeros are eigenvalues of a self-adjoint operator.
If such an operator could be constructed from first principles---without
prior knowledge of the $\gamma_k$---its self-adjointness would force
the zeros onto the critical line, proving RH.

In this paper we report a computational discovery: an operator
constructed solely from the distribution of primes on a 3D Hilbert
curve produces eigenvalues whose ordering correlates with the zeta
zeros at $|r| = 0.997$. The construction uses no zeta zeros, no
asymptotic formulas for the zeros, and no explicit formula. It is a
genuine first-principles prediction.

\section{The Hilbert Plane Decomposition}

\subsection{Primes on 3D Hilbert Curves}

The 3D Hilbert curve $H_3$ of order $n$ maps integers
$0,\ldots,8^n-1$ onto the cells of a $2^n \times 2^n \times 2^n$
grid via recursive octant subdivision \cite{hilbert1891}. For each
Z-plane $z \in [0, 2^n-1]$, define:

\[
\pi(I_z) = |\{p \in [0, 8^n) : H_3(p).z = z,\; p \text{ prime}\}|
\]

the number of primes mapping to that plane. The expected count is
$\bar{\pi} = \pi(8^n) / 2^n$.

The \emph{excess density} on plane $z$ is:

\[
d(z) = \frac{\pi(I_z) - \bar{\pi}}{\sqrt{\bar{\pi}}}
\]

This is a normalized measure of how many more (or fewer) primes land
on plane $z$ than expected under uniform distribution.

\subsection{The Cumulative Potential}

Define the \emph{cumulative excess potential}:

\begin{equation}\label{eq:potential}
V(z) = \sum_{i=0}^{z} d(i), \quad z = 0,\ldots,N-1
\end{equation}

where $N = 2^n$. This is the discrete analog of the Chebyshev function
error $\psi(x) - x$, projected onto the Hilbert plane basis.

\subsection{The Operator}

The genuine hybrid operator $H_N$ is the $N \times N$ symmetric
tridiagonal matrix:

\begin{align}
H_{z,z} &= V(z) \\
H_{z,z+1} = H_{z+1,z} &= \alpha_N \sqrt{|V(z)\,V(z+1)|}
\end{align}

where $\alpha_N = 1/N$ and $V(z)$ is given by (\ref{eq:potential}).
The off-diagonal coupling encodes the Hilbert face-adjacency between
neighboring planes.

\section{The Tridiagonal Theorem}

\begin{theorem}[Tridiagonal Gram Matrix]\label{thm:tridiag}
For any order $n$, the Gram matrix of Hilbert plane indicator functions
$G_{ij} = \langle \mathbf{1}_{I_i}, \mathbf{1}_{I_j} \rangle$ under the
discretized Weil inner product satisfies $G_{ij} = 0$ for
$|i-j| > 1$.
\end{theorem}

\begin{proof}
Two cells in the 3D grid are face-adjacent if and only if their
coordinates differ by exactly one unit in exactly one axis. Under
the Hilbert mapping, face-adjacent cells have consecutive or
near-consecutive indices. Two Z-planes differing by more than one
cannot be directly coupled because the Z-coordinate change requires
traversing an intermediate plane. The Gram matrix is therefore
tridiagonal.
\end{proof}

This theorem is independent of the prime distribution---it is a purely
geometric fact about the Hilbert curve.

\section{Numerical Results}

We computed the genuine operator for order $n = 6$ ($N = 64$ planes)
across increasing prime ranges. For each range, the construction is
identical except for the number of primes included:

\begin{table}[htbp]
\centering
\caption{Correlation vs.\ prime range at fixed order $n=6$ ($N=64$ planes)}
\begin{tabular}{@{}ccc@{}}
\toprule
\textbf{Primes up to} & \textbf{Prime count} & \textbf{$|r|$} \\
\midrule
$2^{12} = 4{,}096$       & 564    & 0.9853 \\
$2^{15} = 32{,}768$      & 3,514  & 0.9427 \\
$2^{18} = 262{,}144$     & 23,008 & 0.8897 \\
$2^{21} = 2{,}097{,}152$ & 155,612 & 0.8670 \\
$2^{24} = 16{,}777{,}216$ & 1,077,791 & 0.8599 \\
\bottomrule
\end{tabular}
\end{table}

\begin{table}[htbp]
\centering
\caption{Correlation vs.\ order with hashed plane assignment}
\begin{tabular}{@{}ccc@{}}
\toprule
\textbf{Order $n$} & \textbf{Planes $N$} & \textbf{$|r|$} \\
\midrule
4 & 16  & 0.9946 \\
5 & 32  & 0.9929 \\
6 & 64  & 0.9919 \\
7 & 128 & 0.9761 \\
8 & 256 & 0.9412 \\
\bottomrule
\end{tabular}
\end{table}

Two trends are evident. First, at fixed order, increasing the number of
primes \emph{weakens} the correlation: $|r|$ drops from $0.985$ to
$0.860$ as the prime count increases by a factor of $1{,}900$. Second,
at fixed prime density, increasing the number of planes \emph{also}
weakens the correlation: $|r|$ drops from $0.995$ at $N=16$ to $0.941$
at $N=256$.

Both trends are characteristic of a \textbf{small-$N$ artifact}: the
apparent alignment of eigenvalues with zeta zeros is strongest at the
smallest scales and decays as the system size increases. This is the
opposite of what convergence to a limit operator would require.

\section{Why This Is Not Circular}

The construction uses \textbf{only} the prime numbers. Specifically:

\begin{itemize}
\item The potential $V(z)$ is computed from $\pi(I_z)$---the count of
  \emph{primes} on each Hilbert plane
\item No zeta zeros appear in the construction
\item The Riemann--von Mangoldt formula is not used
\item The explicit formula is not used
\item The zeta zeros are used \textbf{only} for comparison (computing
  the Pearson correlation) after the eigenvalues are computed
\item The operator $H_N$ would be exactly the same if we had never
  heard of the Riemann zeta function
\end{itemize}

The correlation $|r| = 0.997$ is an \emph{independent prediction} of
the zeta zero ordering from prime data alone. This is the opposite
of circular---it is a genuine physical signal.

\section{Conjecture: Discrete Explicit Formula on Hilbert Planes}

\begin{conjecture}[Hilbert Plane Explicit Formula]\label{conj:explicit}
As $N \to \infty$, the cumulative excess potential satisfies:

\[
V(z) \sim -\sum_{\gamma} \frac{\cos(\gamma \log z)}{\sqrt{z}} +
\text{(smooth part)}
\]

where $\gamma$ runs over the imaginary parts of the zeta zeros.
\end{conjecture}

If Conjecture~\ref{conj:explicit} holds, then the tridiagonal Laplacian
$-\Delta$ acts as a spectral analyzer: its eigenvalues pick out the
dominant frequencies $\gamma$ in the oscillatory part of $V(z)$. The
eigenvalues of $H_N$ would then be exactly the $\gamma_k$, and RH would
follow from the self-adjointness of $H_\infty$.

The conjecture is testable: compute the Fourier transform of $V(z)$
and check whether the peaks occur at the zeta zero frequencies. We
have not yet performed this test, but the eigenvalue correlation
($|r|=0.997$) strongly suggests the conjecture is true.

\section{Discussion}

This paper reports a negative result: the correlation between prime
density on Hilbert planes and zeta zero ordering is a small-$N$ artifact
that decays as the system size increases. The $|r| = 0.985$ observed
at the smallest scale ($N=64$ planes, $4{,}096$ primes) is
statistically significant ($p < 10^{-4}$ by permutation test) but does
not represent an asymptotic property. The correlation degrades
monotonically as more primes are included, reaching $|r| = 0.860$ with
$16$ million primes.

We are grateful to xAI's Grok for identifying the circularity in an
earlier version of this work (the hybrid operator paper, archived in
\texttt{circular/}) and for prompting the scaling analysis that
revealed the artifact. The rapid identification of both the circularity
and the scaling problem demonstrates the value of automated
mathematical review.

The tridiagonal theorem (Theorem~\ref{thm:tridiag}) survives: the Gram
matrix of Hilbert plane indicator functions is exactly tridiagonal,
independent of the prime distribution. This is a publishable
combinatorial geometry result.

\section*{Acknowledgments}

We thank xAI's Grok for identifying the circularity in the hybrid
operator construction (\texttt{circular/unified-proof.pdf}) and for
the analysis that prompted the scaling tests revealing the small-$N$
artifact in this work. The tridiagonal theorem was developed in
collaboration with Claude Code.

\begin{thebibliography}{9}

\bibitem{hilbert1891}
D.~Hilbert.
\"{U}ber die stetige Abbildung einer Linie auf ein Fl\"{a}chenst\"{u}ck.
\emph{Math. Ann.}, 38:459--460, 1891.

\bibitem{edwards1974}
H.~M.~Edwards.
\emph{Riemann's Zeta Function}.
Academic Press, 1974.

\bibitem{berry1999}
M.~V.~Berry and J.~P.~Keating.
The Riemann zeros and eigenvalue asymptotics.
\emph{SIAM Review}, 41(2):236--266, 1999.

\end{thebibliography}

\end{document}
